% GENERATED FILE. Edit canonical lessons, then run npm run generate:books.
% canonical-source-hash: fnv1a64:e5d5bd526ba23e68
% canonical-lessons: MR-C497-dhakta, MR-C497-thorla, MR-C497-ekulta, MR-C497-culat, MR-C497-savatra

\chapter{Elder and Younger}
\label{ch:mr-elder-and-younger}
\hlchaptermodality{497}

\begin{chapteropening}
\textbf{By the end of this chapter:} \emph{I can say who is younger or elder, an only child, a cousin or a step-relation.}

Complete the last lesson of chapter 497: I can say who is younger or elder, an only child, a cousin or a step-relation.
\end{chapteropening}

\section[\texorpdfstring{\mr{धाकटा} (dhākṭā)}{धाकटा (dhākṭā)}]{\mr{धाकटा} (dhākṭā) — younger}

\label{lesson:MR-C497-dhakta}



\begin{quote}
Before the new one: say the Marathi for a brother's wife, then the Marathi for a child's father-in-law.
\end{quote}

\subsection*{You'll want to know: \mr{धाकटा}}
\textbf{\mr{धाकटा}} — \emph{dhākṭā} — \textquotedblleft{}younger\textquotedblright{}.

\begin{quote}
\mr{माझा} \mr{धाकटा} \mr{भाऊ} \mr{शाळेत} \mr{जातो}.
\end{quote}

\begin{quote}
\emph{mājhā dhākṭā bhāū śāḷet jāto.}
\end{quote}

\begin{quote}
My younger brother goes to school.
\end{quote}

Agrees: \mr{धाकटा} \mr{भाऊ}, \mr{धाकटी} \mr{बहीण}.

\begin{cousinweb}
\textbf{\mr{धाकटा}} (\emph{dhākṭā}) is a Marathi adjective whose origin is not settled.
\end{cousinweb}

\begin{grammarlens}[title={What you've built}]
One more word for describing what you lived and felt.
\end{grammarlens}

\subsection*{Your turn}
\begin{itemize}
\raggedright
  \item \emph{Say it:} \emph{dhākṭā}
  \item \emph{Say it:} \emph{dhākṭā}, once more
  \item \emph{Say it:} \emph{mājhā dhākṭā bhāū śāḷet jāto.} — My younger brother goes to school.
  \item \emph{Recall:} say the Marathi for a brother's wife, then the Marathi for a child's father-in-law, then say \emph{dhākṭā} again
\end{itemize}

\begin{tcolorbox}[breakable,colback=teal!4,colframe=teal!35!black,title={Before you move on}]
What does \mr{धाकटा} mean? (\textquotedblleft{}Younger\textquotedblright{}.) Say it once more.
\end{tcolorbox}
\section[\texorpdfstring{\mr{थोरला} (thorlā)}{थोरला (thorlā)}]{\mr{थोरला} (thorlā) — elder}

\label{lesson:MR-C497-thorla}



\begin{quote}
Before the new one: say the Marathi for a child's father-in-law, then the Marathi for younger.
\end{quote}

\subsection*{You'll want to know: \mr{थोरला}}
\textbf{\mr{थोरला}} — \emph{thorlā} — \textquotedblleft{}elder\textquotedblright{}.

\begin{quote}
\mr{थोरल्या} \mr{बहिणीने} \mr{मला} \mr{गणित} \mr{शिकवलं}.
\end{quote}

\begin{quote}
\emph{thorlyā bahiṇīne malā gaṇit śikavlaṁ.}
\end{quote}

\begin{quote}
My elder sister taught me maths.
\end{quote}

Agrees; its pair is \mr{धाकटा}.

\begin{cousinweb}
\textbf{\mr{थोरला}} (\emph{thorlā}) is a Marathi adjective whose origin is not settled.
\end{cousinweb}

\begin{grammarlens}[title={What you've built}]
One more word for describing what you lived and felt.
\end{grammarlens}

\subsection*{Your turn}
\begin{itemize}
\raggedright
  \item \emph{Say it:} \emph{thorlā}
  \item \emph{Say it:} \emph{thorlā}, once more
  \item \emph{Say it:} \emph{thorlyā bahiṇīne malā gaṇit śikavlaṁ.} — My elder sister taught me maths.
  \item \emph{Recall:} say the Marathi for a child's father-in-law, then the Marathi for younger, then say \emph{thorlā} again
\end{itemize}

\begin{tcolorbox}[breakable,colback=teal!4,colframe=teal!35!black,title={Before you move on}]
What does \mr{थोरला} mean? (\textquotedblleft{}Elder\textquotedblright{}.) Say it once more.
\end{tcolorbox}
\section[\texorpdfstring{\mr{एकुलता} (ekultā)}{एकुलता (ekultā)}]{\mr{एकुलता} (ekultā) — the only one}

\label{lesson:MR-C497-ekulta}



\begin{quote}
Before the new one: say the Marathi for younger, then the Marathi for elder.
\end{quote}

\subsection*{You'll want to know: \mr{एकुलता}}
\textbf{\mr{एकुलता}} — \emph{ekultā} — \textquotedblleft{}the only one\textquotedblright{}.

\begin{quote}
\mr{मी} \mr{माझ्या} \mr{आईवडिलांचा} \mr{एकुलता} \mr{एक} \mr{मुलगा} \mr{आहे}.
\end{quote}

\begin{quote}
\emph{mī mājhyā āīvaḍilāñcā ekultā ek mulgā āhe.}
\end{quote}

\begin{quote}
I am my parents' only son.
\end{quote}

With \mr{एक}: \mr{एकुलता} \mr{एक} \mr{मुलगा}, \mr{एकुलती} \mr{एक} \mr{मुलगी}.

\begin{cousinweb}
\textbf{\mr{एकुलता}} (\emph{ekultā}) is a Marathi adjective whose origin is not settled.
\end{cousinweb}

\begin{grammarlens}[title={What you've built}]
One more word for describing what you lived and felt.
\end{grammarlens}

\subsection*{Your turn}
\begin{itemize}
\raggedright
  \item \emph{Say it:} \emph{ekultā}
  \item \emph{Say it:} \emph{ekultā}, once more
  \item \emph{Say it:} \emph{mī mājhyā āīvaḍilāñcā ekultā ek mulgā āhe.} — I am my parents' only son.
  \item \emph{Recall:} say the Marathi for younger, then the Marathi for elder, then say \emph{ekultā} again
\end{itemize}

\begin{tcolorbox}[breakable,colback=teal!4,colframe=teal!35!black,title={Before you move on}]
What does \mr{एकुलता} mean? (\textquotedblleft{}The only one\textquotedblright{}.) Say it once more.
\end{tcolorbox}
\section[\texorpdfstring{\mr{चुलत} (culat)}{चुलत (culat)}]{\mr{चुलत} (culat) — of a paternal uncle; cousin on the father's side}

\label{lesson:MR-C497-culat}



\begin{quote}
Before the new one: say the Marathi for elder, then the Marathi for the only one.
\end{quote}

\subsection*{You'll want to know: \mr{चुलत}}
\textbf{\mr{चुलत}} — \emph{culat} — \textquotedblleft{}of a paternal uncle; cousin on the father's side\textquotedblright{}.

\begin{quote}
\mr{माझा} \mr{चुलत} \mr{भाऊ} \mr{आणि} \mr{मी} \mr{एकाच} \mr{वर्गात} \mr{आहोत}.
\end{quote}

\begin{quote}
\emph{mājhā culat bhāū āṇi mī ekāc vargāt āhot.}
\end{quote}

\begin{quote}
My cousin and I are in the same class.
\end{quote}

Unchanging: \mr{चुलत} \mr{भाऊ}, \mr{चुलत} \mr{बहीण}, the children of a father's brother.

\begin{cousinweb}
\textbf{\mr{चुलत}} (\emph{culat}) is inherited: CDIAL lists Marathi \emph{cultā}, \textquotedblleft{}father's brother\textquotedblright{}, under the reconstructed \emph{*culla}, \textquotedblleft{}small\textquotedblright{} (CDIAL 4877).
\end{cousinweb}

\begin{grammarlens}[title={What you've built}]
One more word for describing what you lived and felt.
\end{grammarlens}

\subsection*{Your turn}
\begin{itemize}
\raggedright
  \item \emph{Say it:} \emph{culat}
  \item \emph{Say it:} \emph{culat}, once more
  \item \emph{Say it:} \emph{mājhā culat bhāū āṇi mī ekāc vargāt āhot.} — My cousin and I are in the same class.
  \item \emph{Recall:} say the Marathi for elder, then the Marathi for the only one, then say \emph{culat} again
\end{itemize}

\begin{tcolorbox}[breakable,colback=teal!4,colframe=teal!35!black,title={Before you move on}]
What does \mr{चुलत} mean? (\textquotedblleft{}Of a paternal uncle; cousin on the father's side\textquotedblright{}.) Say it once more.
\end{tcolorbox}
\section[\texorpdfstring{\mr{सावत्र} (sāvatra)}{सावत्र (sāvatra)}]{\mr{सावत्र} (sāvatra) — step-}

\label{lesson:MR-C497-savatra}



\begin{quote}
Before the new one: say the Marathi for the only one, then the Marathi for of a paternal uncle.
\end{quote}

\subsection*{You'll want to know: \mr{सावत्र}}
\textbf{\mr{सावत्र}} — \emph{sāvatra} — "step-".

\begin{quote}
\mr{त्याची} \mr{सावत्र} \mr{आई} \mr{त्याच्यावर} \mr{खूप} \mr{प्रेम} \mr{करते}.
\end{quote}

\begin{quote}
\emph{tyācī sāvatra āī tyācyāvar khūp prem karte.}
\end{quote}

\begin{quote}
His stepmother loves him very much.
\end{quote}

Unchanging: \mr{सावत्र} \mr{आई}, a stepmother; \mr{सावत्र} \mr{भाऊ}, a half-brother.

\begin{cousinweb}
\textbf{\mr{सावत्र}} (\emph{sāvatra}) is a tadbhava: CDIAL lists Marathi \emph{sāvatra}, \textquotedblleft{}born of a co-wife\textquotedblright{}, under Sanskrit \emph{sāpatna}, \textquotedblleft{}of a rival\textquotedblright{} (CDIAL 13342).
\end{cousinweb}

\begin{grammarlens}[title={What you've built}]
One more word for describing what you lived and felt.
\end{grammarlens}

\subsection*{Your turn}
\begin{itemize}
\raggedright
  \item \emph{Say it:} \emph{sāvatra}
  \item \emph{Say it:} \emph{sāvatra}, once more
  \item \emph{Say it:} \emph{tyācī sāvatra āī tyācyāvar khūp prem karte.} — His stepmother loves him very much.
  \item \emph{Recall:} say the Marathi for the only one, then the Marathi for of a paternal uncle, then say \emph{sāvatra} again
\end{itemize}

\begin{tcolorbox}[breakable,colback=teal!4,colframe=teal!35!black,title={Before you move on}]
What does \mr{सावत्र} mean? ("Step-".) Say it once more.
\end{tcolorbox}
